\statement{A.137}{Proved}{The original X-direction negative frames and the exact modular nonextension}
Retain the ORIGINAL negative full-\(\widehat{GT}\) involution,
its actual elliptic action \(\Theta\), and the two based source maps
of A.134--A.135:
\[
 T=\widehat{(B_4/Z(B_4))\times_{S_4}S_3},
 \qquad \jmath:T\longrightarrow G=\widehat B_{1,3},
 \qquad q:T\longrightarrow B=C_2(s)\amalg C_3(r).
\]
The actual X-direction cubic has full normalizer \(S_3\) in
\(G\rtimes_\Theta C_2\). Its correctly normalized negative frame
and the ORIGINAL modular negative frame are the two designated
images of one source \(S_3\).
Nevertheless, for every profinite group \(H\) and continuous
\(\rho:F=\ker(B\to S_3)\to H\), the pure source pullback
\(\rho q|_{P_T}\) extends across \(\jmath|_{P_T}\) to the whole
elliptic pure group \(P=\ker(G\to S_3)\) if and only if
\(\rho=1\). In particular the original signed \(q_3\) does not
have such an extension.
\begin{proof}
\emph{1. The actual X cubic and its negative normalizer.}
Write
\[
 a=\sigma^2,\quad b=\tau^2,\quad
 X_2=\sigma X\sigma,\quad X_3=\tau X_2\tau,\quad
 R_X=X\sigma\tau .
\]
The original relations are
\[
 XX_2X_3=1,\qquad [X,\tau]=1,\qquad
 \sigma\tau\sigma=\tau\sigma\tau.
\]
They give the complete ordered calculation
\[
\begin{aligned}
 R_X^3
 &=X\sigma\tau X\sigma\tau X\sigma\tau\\
 &=X\sigma X\tau\sigma\tau X\sigma\tau\\
 &=X\sigma X\sigma\tau\sigma X\sigma\tau\\
 &=X(\sigma X\sigma)\tau(\sigma X\sigma)\tau
   =XX_2X_3=1 .
\end{aligned}
\tag{A.137.1}
\]
Its permutation is \(c=(12)(23)\), so its order is exactly three.
A.124 also gives its nonzero abelian torsion label \(X\), distinct
from the label \(Y\) of \(\Pi_0=\tau^{-1}\sigma^{-1}Y\), and gives
\[
 C_G(R_X)=N_G(\langle R_X\rangle)=\langle R_X\rangle,
 \qquad \langle R_X\rangle\cap P=1.
\tag{A.137.2}
\]
These are the established full-completion conclusions.

The original \(\Theta\) acts identically on \(S_3\) and acts on
the odd torsion labels by \((X,Y)\mapsto(-X,Y)\).
Thus \(\Theta(R_X)\) and \(\tau R_X^{-1}\tau^{-1}\) have the
same label and the same permutation \(c\). A.124's complete
torsion classification makes them conjugate. A conjugator has
permutation in the centralizer of \(c\), so right multiplication
by a suitable power of the target cubic makes it pure.
The pure centralizer in (A.137.2) is trivial. Hence there is a
unique \(g_\tau\in P\) such that, with \(J_{X,\tau}=g_\tau\tau\),
\[
 \Theta(R_X)=J_{X,\tau}R_X^{-1}J_{X,\tau}^{-1},
 \qquad \pi(J_{X,\tau})=\pi(\tau).
\tag{A.137.3}
\]
Applying \(\Theta\) twice shows that
\(\Theta(J_{X,\tau})J_{X,\tau}\) centralizes \(R_X\).
Its permutation is trivial, so it is one.

Let \(\vartheta\) denote the external involution for \(\Theta\).
Then
\[
 \kappa_X=J_{X,\tau}^{-1}\vartheta,\qquad
 \kappa_X^2=1,\qquad
 \kappa_XR_X\kappa_X^{-1}=R_X^{-1}.
\]
The resulting subgroup is \(S_3\). Its intersection with \(G\)
is the full normalizer \(\langle R_X\rangle\) from (A.137.2);
any element of the other normalizer coset differs from the
known \(\kappa_X\) by an element of that intersection. Therefore
\[
\begin{aligned}
 N_{G\rtimes_\Theta C_2}(\langle R_X\rangle)
   &=\langle R_X,\kappa_X\rangle\cong S_3,\\
 C_{G\rtimes_\Theta C_2}(R_X)&=\langle R_X\rangle .
\end{aligned}
\tag{A.137.4}
\]
This is the X-direction normalizer; it is not the \(C_6\)
normalizer of the original \(\Theta\)-fixed Y cubic in A.133.

The permutation needed for the modular \(s\)-frame is
\(\pi(\sigma\tau\sigma)\), rather than \(\pi(\tau)\).
The exactly normalized transporter is
\[
 J_{X,s}=J_{X,\tau}R_X .
\tag{A.137.5}
\]
Indeed \(\pi(\tau)c=\pi(\sigma\tau\sigma)\); right multiplication
by \(R_X\) preserves the negative conjugacy, and
\[
 \Theta(J_{X,\tau}R_X)J_{X,\tau}R_X
 =J_{X,\tau}^{-1}
   (J_{X,\tau}R_X^{-1}J_{X,\tau}^{-1})
   J_{X,\tau}R_X=1.
\]
The same pure-centralizer argument gives uniqueness with this
specified permutation.

\par\smallskip\noindent
\emph{2. Both exact source normalizations.}
Use the A.134--A.135 generator table
\[
\begin{array}{c|ccc}
 &t&\sigma&\tau\\\hline
 \jmath&X&\sigma&\tau\\
 q&1&r^{-1}s&sr^{-1}.
\end{array}
\tag{A.137.6}
\]
Put \(R=t\sigma\tau\), so that \(\jmath(R)=R_X\) and
\(q(R)=r\). The source presentation has \(R^3=1\).
The complete source torsion proof in A.135 gives two conjugacy
classes of nonidentity elements of order three, represented
by \(R,R^{-1}\), and
\[
 C_T(R)=N_T(\langle R\rangle)=\langle R\rangle,
 \qquad \langle R\rangle\cap P_T=1,
 \quad P_T=\ker(T\to S_3).
\tag{A.137.7}
\]
The actual involution \(\theta_T\) acts by minus one on the
source abelianization; in particular it sends the nonzero
torsion label of \(R\) to its negative. Consequently
\(\theta_T(R)\) is conjugate to \(R^{-1}\).
Normalize a conjugator to permutation \(\pi(\tau)\), exactly
as in part 1, and call it \(J_\tau\). Then
\[
 \theta_T(R)=J_\tau R^{-1}J_\tau^{-1},\qquad
 \theta_T(J_\tau)J_\tau=1.
\]
The correct second normalization is
\[
 J_s=J_\tau R,\qquad
 \pi(J_s)=\pi(\sigma\tau\sigma),\qquad
 \theta_T(J_s)J_s=1 .
\tag{A.137.8}
\]
These are unique source transporters in their specified
permutation cosets.

The covariance in A.135 is
\[
 \jmath\theta_T=\Theta_0\jmath,\qquad
 q\theta_T=\alpha q,\qquad
 \Theta_0=\operatorname{Ad}(h_1^{-1})\Theta .
\]
Retain the ORIGINAL modular negative transporter
\[
 J_1^\alpha=s\phi A_r,\qquad
 \alpha(r)=J_1^\alpha r^{-1}(J_1^\alpha)^{-1}
\tag{A.137.9}
\]
from A.111, with its original normalization.
Both \(q(J_s)\) and \(J_1^\alpha\) conjugate \(r^{-1}\)
to \(\alpha(r)\), with permutation \(s\).
Their quotient belongs to \(C_B(r)\cap F=1\), where
\(F=\ker(B\to S_3)\). Thus the exact images are
\[
 q(J_s)=J_1^\alpha,\qquad
 q(J_\tau)=J_1^\alpha r^{-1}.
\tag{A.137.10}
\]
The \(R\) on the right of \(J_\tau\) in (A.137.8) is
essential; omitting it changes the marked modular frame.
Similarly \(\jmath(J_s)\) is the unique \(s\)-normalized
negative transporter for \(\Theta_0\).

\par\smallskip\noindent
\emph{3. The ORIGINAL \(\Theta\) comparison, including the
external involution.}
In the source use the actual A.135 words
\[
\begin{gathered}
 K_S=\sigma\tau^2\sigma,\quad
 U_S=u(t,K_S),\quad Z_{u,S}=u(\sigma t\sigma,\tau^2),\\
 r_S=U_SZ_{u,S},\quad
 h_S=r_S\theta_T(r_S)^{-1}
     =\phi(t,K_S)\phi(\sigma t\sigma,\tau^2),\\
 \beta_T=\operatorname{Ad}(h_S)\theta_T .
\end{gathered}
\tag{A.137.11}
\]
These are literal source identities already proved in A.135.
They give
\[
 \jmath\beta_T=\Theta\jmath,\qquad
 q\beta_T=\operatorname{Ad}(\phi)\alpha q,\qquad
 \jmath(h_S)=h_1,\quad q(h_S)=\phi .
\tag{A.137.12}
\]
The original source transporter is
\[
 J_s^{\rm orig}=h_SJ_s .
\]
It has the same permutation as \(J_s\), and direct substitution
gives
\[
 \beta_T(R)=J_s^{\rm orig}R^{-1}(J_s^{\rm orig})^{-1},
 \qquad
 \beta_T(J_s^{\rm orig})J_s^{\rm orig}=1.
\tag{A.137.13}
\]
For example the second equality uses
\(h_S\theta_T(h_S)=1\) and \(\theta_T(J_s)J_s=1\).
Uniqueness in (A.137.3)--(A.137.5) identifies
\(\jmath(J_s^{\rm orig})=J_{X,s}\).
Its modular image is
\[
 q(J_s^{\rm orig})=\phi J_1^\alpha=sA_r.
\tag{A.137.14}
\]
Here \(\phi s\phi=s\) follows from the original
\(\alpha(s)=s\phi\) and \((s\phi)^2=1\).

Let \(\vartheta_S\) be the external involution for \(\beta_T\),
and \(\vartheta_\alpha\) the external involution for \(\alpha\).
There are genuine continuous homomorphisms
\[
\begin{array}{rcl}
 T\rtimes_{\beta_T} C_2&\longrightarrow&
          G\rtimes_\Theta C_2,\\
 (t_0,\vartheta_S)&\longmapsto&
          (\jmath(t_0),\vartheta),
\end{array}
\qquad
\begin{array}{rcl}
 T\rtimes_{\beta_T} C_2&\longrightarrow&
          B\rtimes_\alpha C_2,\\
 (t_0,\vartheta_S)&\longmapsto&
          (q(t_0),\phi\vartheta_\alpha).
\end{array}
\tag{A.137.15}
\]
This notation specifies the images of every source element
\(t_0\) and of the external generator. The covariance is
(A.137.12); the second external generator has square one because
\(\phi\alpha(\phi)=1\), by the original canonical norm.
The specified source reflection has images
\[
 (J_s^{\rm orig})^{-1}\vartheta_S
 \longmapsto J_{X,s}^{-1}\vartheta,\qquad
 (J_s^{\rm orig})^{-1}\vartheta_S
 \longmapsto (J_1^\alpha)^{-1}\vartheta_\alpha .
\tag{A.137.16}
\]
The latter equality retains the cancellation
\((\phi J_1^\alpha)^{-1}\phi=(J_1^\alpha)^{-1}\).
The source subgroup generated by \(R\) and this reflection is
\(S_3\), by (A.137.13). Both maps restrict to isomorphisms onto
the indicated actual negative frames: the cubic stays nontrivial
and the reflection has nontrivial external coordinate. This is a
comparison through the actual common source, not a map from the
whole elliptic group to the modular group.

\par\smallskip\noindent
\emph{4. The precise obstruction to extending a modular pure
quotient.}
The restrictions of the two maps in (A.137.6) have domains
\[
 \jmath:P_T\longrightarrow P,\qquad
 q:P_T\longrightarrow F.
\]
The latter is onto, since \(a,b\) map to the original free
generators
\[
 x=(r^{-1}s)^2,\qquad y=(sr^{-1})^2 .
\]
For \(t_2=\sigma t\sigma\), the required original images are
\[
\begin{array}{c|cc}
 \text{source word}&\jmath&q\\\hline
 t&X&1\\
 t_2&X_2&x\\
 a&a&x\\
 K_S&K=[Y,X]&y^{-1}.
\end{array}
\tag{A.137.17}
\]
The last row follows from the three-strand braid identity
\((\sigma\tau)^3=bK_S\); its modular image is one and
\(q(b)=y\). It does not use centrality of that braid word
in the whole elliptic group.

Suppose that a continuous \(\eta:P\to H\) satisfies
\[
 \eta\jmath|_{P_T}=\rho q|_{P_T}.
\tag{A.137.18}
\]
The literal PURE elliptic relations from A.121 are
\[
 K=[Y,X],\qquad [Y,X_2^{-1}]=a .
\tag{A.137.19}
\]
The first and last rows of (A.137.17) give
\[
 \eta(X)=1,\qquad
 \rho(y)^{-1}=\eta(K)=[\eta(Y),1]=1.
\]
The middle rows give
\(\eta(X_2)=\eta(a)=\rho(x)=M\).
Writing \(V=\eta(Y)\), the second relation in (A.137.19)
becomes
\[
 [V,M^{-1}]=M,\qquad VM^{-1}V^{-1}M=M .
\]
Right cancellation gives \(VM^{-1}V^{-1}=1\), hence \(M=1\).
Both free-generator images of \(\rho\) are therefore one,
so \(\rho=1\). Conversely the trivial map \(P\to H\)
extends the trivial \(\rho\). This proves the stated
if-and-only-if criterion for arbitrary profinite \(H\).

For the original signed representation of A.105/A.111,
\[
 q_3(x)=\begin{pmatrix}1&2\\0&1\end{pmatrix},
 \qquad
 q_3(y)=\begin{pmatrix}1&0\\1&1\end{pmatrix}
 \quad\text{in }\mathrm{SL}_2(\mathbf F_3).
\tag{A.137.20}
\]
These images are nontrivial, so its pure source pullback has
no extension to \(P\), and therefore none to \(G\).
The notation uses \(q_3\) on \(F\) only; it does not define
an unproved \(q_3\) on \(B\) or on all of \(T\).
The same obstruction holds for a conjugate of this marked
representation by conjugating a hypothetical extension back.
\end{proof}

\paragraph{Source consolidation and remaining scope.}
The distinct conclusions consolidate the accepted sources
\path{work_A/actual_elliptic_X_cubic_negative_frame37.md}
(\texttt{75f055e0}),
\path{work_DE/actual_annular_X_frame_modular_comparison37.md}
(\texttt{7f758ffb}),
\path{work_DE/actual_annular_original_Theta_frame_upgrade37.md}
(\texttt{a3c81a99}), and
\path{work_A/actual_annular_forgetting_pure_extension_obstruction37.md}
(\texttt{a0b3ef4f}). The source action, topology, torsion and gauge
proofs already printed in A.134--A.135 are reused at their exact
scopes. The comparison transports the existing negative frame;
it supplies no canonical positive frame or original CM sign.
The nonextension criterion concerns only the stated marked
source restrictions, not every unrelated representation or
possible sewing construction. Main A remains Open.
